\documentclass[../main.tex]{subfiles}
\begin{document}
%==========================================TIÊU ĐỀ TẬP========================================
\begin{table}[h]
  \begin{adjustwidth}{-1cm}{}
    \begin{tabular}{>{\centering\arraybackslash}p{6cm}|>{\raggedright\arraybackslash}p{12.5cm}}
      \multirow{5}{6cm}
      % Cột bên trái: ÉPISODE và số tập
      {\\[-40pt]
        \flaregothic\fontsize{20pt}{20pt}\selectfont \centering
        \color{eptype}ÉPISODE\color{black}\\
        \burbank\fontsize{72pt}{72pt}\selectfont
        \color{epnum}74\color{black}\\[6pt]
        \cafeteria\fontsize{20pt}{20pt}\selectfont\color{epnum}(6-5)\color{black}
        \\[18pt]
        % \flaregothic\fontsize{14pt}{14pt}\selectfont
        % \color{epattr}BIENVENUE, \uline{\color{Banpire} Mel} !
        % \flaregothic\fontsize{18pt}{18pt}\selectfont
        % \color{epattr}ÉPISODE PERDU
        \color{black}
      }
       &
      % Dòng thứ nhất: tên tập tiếng Nhật
      {\fotskip\fontsize{18pt}{18pt}\selectfont バンドやろうぜ！
      }                                                        \\[6pt]
      % Dòng thứ hai: tên tập tiếng Anh
       & {\cafeteria\fontsize{18pt}{18pt}\selectfont Triangle}                                                 \\[4pt]
      % Dòng thứ ba: tên tập tiếng Việt
       & {\vnmsans\fontsize{18pt}{18pt}\selectfont Kẻng}                                                       \\[8pt]
      % Dòng thứ tư: Nhân vật xuất hiện
       & {\caviar{\fontsize{14pt}{14pt}\selectfont Nhân vật: \emp{kuon1} \emp{ryuuhou1} \emp{aki} \emp{matsuri} \emp{mio} \emp{rushia}}
         }
      \\[8pt]
      % Dòng cuối cùng: Thời gian diễn ra của tập (thường sẽ là ngày phát hành)
       & {\continuum\fontsize{16pt}{16pt}\selectfont Ngày phát hành: 04/10/2020}                               \\[18pt]
    \end{tabular}
  \end{adjustwidth}
\end{table}
%=========================================TRUYỆN CHÍNH========================================
\color{black}

\txtbx[2]
{{\color{vol} \uline{Tóm tắt:}} Matsuri bỗng dưng muốn lập ban nhạc với một cái kẻng thần kỳ, nhưng tiếng đầu tiên lại phát ra ``bụp'' thảm hại. Sau một hồi thiền tịnh ``thanh lọc tâm hồn'' dưới sự hướng dẫn của Mio, cô gõ ra tiếng ``KENGGGG'' vang động cả Trái Đất. Thế là nhóm ``Chị em Hút bụi Ion âm'' ra đời với Aki (vĩ cầm), Mio (maraca) và Matsuri (kẻng). Rushia không chịu được, hét vào cùng, còn Ryuuhou - trưởng nhóm underground có kinh nghiệm - cầm ukulele đệm vào cho đỡ chán. Kết quả: một buổi tập nhạc mà chẳng ai dám gọi là nhạc, nhưng vui hơn cả một buổi hòa nhạc thật.
}

Một ngày đầu tháng Mười. Những cơn gió đầu thu bắt đầu mang theo cái se lạnh, khiến lá vàng rơi đầy sân sau văn phòng, từng chiếc lá nhẹ nhàng đáp xuống như những nốt nhạc lặng của bản giao hưởng mùa thu. Trời hôm nay không quá nóng cũng không quá lạnh, một thời điểm lý tưởng để ngồi nhâm nhi tách trà nóng hoặc đọc báo, hoặc đơn giản là ngồi không mà chẳng làm gì cả.

Căn phòng lúc này yên tĩnh đến mức chỉ nghe thấy tiếng bộp bộp của màn hình điện thoại và quân cờ va vào mặt bàn, một bản hòa tấu nhẹ nhàng của sự thư thái.

Fubuki đang lướt điện thoại, Mio và Rushia ngồi đối diện nhau trên bàn cờ vây, còn cậu Tổng thì đang cầm tờ báo, nhìn chằm chằm vào trang thể thao như thể đó là bí mật của vũ trụ. Ryuuhou ngồi bên cạnh Kuon, đang vẽ nguệch ngoạc gì đó lên một mảnh giấy, trông như một con thỏ có cánh hoặc một con rồng đang đội mũ, nhưng chẳng ai đủ can đảm để hỏi cô đang vẽ cái gì. Thỉnh thoảng cô ngước lên nhìn anh trai, khẽ cười, rồi lại cúi xuống tiếp tục tác phẩm của mình. Hai tai của cô thì bịt tai nghe, phát những bản nhạc từ nhóm nhạc underground của mình.

Bất ngờ, Matsuri từ phía bên ngoài bước ngoài bỗng reo lên như vừa nghĩ ra một sáng kiến: ``\textbf{Mình cùng lập một ban nhạc nào!}''

Nàng Sói đen và nàng Pháp sư đang ngồi đánh cờ vây, còn cậu Tổng đang đọc báo. Nghe cô nói vậy, cả ba người đồng loạt ngẩng lên, nhìn cô với vẻ mặt hững hờ, không nói một lời nào, như thể họ vừa nghe thấy một câu đùa quá cũ. Không khí im lặng kéo dài trong vài giây, đủ dài để một con ốc sên bò qua bàn cờ.

Cô nàng mắt hai màu dừng bút, ngước nhìn lên, đôi tai vểnh lên vì tò mò: ``Ban nhạc ạ? Matsuri-san định hát gì thế?'' với vẻ mặt đầy hào hứng, hoàn toàn trái ngược với sự thờ ơ của ba người kia.

Không khí im lặng kéo dài trong vài giây, như thể cả căn phòng đang cân nhắc xem câu nói đó có đáng để phản hồi hay không.

\img{6-5a}

Nàng Lễ hội hơi nhíu mày rồi đặt tay lên vai hậu bối, hạ giọng nói nhẹ nhàng hơn, như thể cô là nhân vật chính của bộ phim tình cảm lãng mạn đang chiếu giữa khung cảnh mùa thu: ``Hãy lập một ban nhạc với chị đi.''

\img{6-5b}

Cậu con trai phì cười, ngước khỏi màn hình điện thoại: ``Khác gì nhau đâu em?''

Trong khi đó, Mio cũng nheo mắt, đôi tai hơi dựng lên: ``Em hiểu ý chị rồi\dots\ Nhưng sao tự dưng lại muốn lập một ban nhạc thế?''

Rushia khoanh tay, gật gù, tỏ vẻ lo ngại: ``Dám cá là từ TV anime mà ra.'' Cậu Tổng bình thản đáp, mắt vẫn dán vào tờ báo: ``Dễ là bắt chước từ \emph{L*ve L*ve!} lắm.''

Matsuri lắc đầu nguầy nguậy, tay đập bàn phủ nhận: ``Không phải! Hôm qua chị xem một chương trình trên Pr*me!'' Cả nàng Sói đen và cậu Tổng đồng thanh thốt lên, giọng đều đều như hai cỗ máy: ``Có khác quái gì nhau đâu?''

Ngay lúc đó, cô em gái tháo một bên tai nghe xuống, cất giọng tinh quái: ``Em thì nghĩ ban nhạc cần có cái tên hay trước, rồi mới tới là có nguồn cảm hứng từ đâu.''

Matsuri quay phắt sang Ryuuhou, mắt sáng rực: ``Ồ! Có cao thủ âm nhạc ở đây này! Em giúp chị nhé!''

Ryuuhou cười nhẹ, đặt tách trà xuống: ``Chị cứ nói ý tưởng của chị đã.''

Bỏ ngoài tai những lời bình luận, nàng Cổ động vẫn cười toe toét và vỗ mạnh vào tờ giấy trên bàn, nhiệt tình không kém gì một diễn viên quảng cáo: ``Dù sao thì chị cũng đã nghĩ ra tên của ban nhạc rồi đây này!''

Nàng Pháp sư liếc nhìn, khoanh tay đoán mò: ``Chắc lại kiểu `Matsuristreet Girls' cho mà xem.''

Mio và Kuon cũng tò mò ghé mắt nhìn tờ giấy, nhưng ngay sau đó cả ba người đều sững sờ. Ryuuhou cũng nghiêng người nhìn, và một tia hoài nghi hiện rõ trên gương mặt cô.

\img{6-5c}

``Gần đúng rồi! Đây này!'' tên mà Matsuri định đặt lại là\dots\ ``Chị em Hút bụi ion âm (マイナスイオン吸引大姉妹)''!

Nàng Sói đen nhăn mặt, đôi mắt híp lại đầy khó hiểu: ``Nghe chả ổn tí nào cả!''

Thậm chí đến cả cậu Tổng, người thường không quan tâm đến mấy trò vớ vẩn, cũng không thể chấp nhận cái tên đó, nhíu mày lắc đầu: ``Sao anh nghe nó chả ăn nhập gì thế\dots''

Ryuuhou khúc khích cười, nhấp một ngụm trà: ``Matsuri-san ơi, cái tên đó nghe như một câu thần chú trong phim bảo vệ môi trường hơn là tên ban nhạc đấy.''

Nhưng như thường lệ, cô nàng Lễ hội chẳng mảy may bận tâm đến phản ứng của mọi người, thậm chí còn hào hứng vung tay chỉ vào một góc phòng: ``Không sao! Chị còn chuẩn bị sẵn cả nhạc cụ nữa đây!''

Cô hớn hở lôi ra một thanh kim loại hình tam giác đơn độc, trông như thể vừa được vớt ra từ một bộ dụng cụ cứu hỏa.

\img{6-5d}

Mio chớp mắt, giọng đầy hoài nghi: ``Ồ, một cái kẻng.'' Kuon khoanh tay, tỏ vẻ hoài nghi hơn nữa: ``Mỗi cái kẻng này thì làm được gì?''

Rushia nhìn Matsuri bằng ánh mắt thắc mắc, như thể đang hỏi một câu hỏi hiển nhiên: ``Cái đó đủ để thành lập một ban nhạc không vậy?''

Ryuuhou gật gù, với kinh nghiệm của một người từng quản lý nhạc cụ thật: ``Một ban nhạc cần ít nhất có trống, guitar, bass, hoặc keyboard chứ. Mỗi cái kẻng thì thành `nhóm kẻng' chứ không phải ban nhạc đâu, chị ơi.''

Nhưng cô nàng Cổ động nhí nhảnh vẫn chẳng có vẻ gì là chùn bước. Cô hắng giọng, vẻ tự tin như thể đã lên kế hoạch cả đời: ``Chị sẽ chơi kẻng! Hai đứa kia thì\dots''

Cô dừng lại, khoanh tay suy nghĩ.

Nàng Pháp sư lẩm bẩm, ngón tay gõ nhẹ lên cằm: ``Có thể em sẽ làm phần giọng hát nhỉ\dots'' Nàng Sói đen nhún vai, đưa mắt nhìn cậu Tổng: ``Chắc anh cũng phải có vai trò gì đó trong nhóm chứ nhỉ?''

Cậu con trai chán nản đáp, vẫn giữ nguyên vẻ mặt thờ ơ: ``Anh chắc là làm quản lý ban nhạc đây\dots''

Cô em gái của cậu nghe vậy, bỗng bật cười nhẹ, giọng đầy ẩn ý: ``Anh từng làm quản lý cho nhóm em, nhớ không? Nếu anh cầm cái kẻng thì e rằng nhóm sẽ tan rã trước khi ra bài đầu tiên đấy.''

Cậu quay sang em gái, nhướng mày: ``Em đang bênh mấy con thú này đấy à?''

``Em chỉ nói sự thật thôi,'' Ryuuhou nhún vai, nụ cười ranh mãnh.

Ngay lúc đó, Nàng Cổ động lôi ra một cây vĩ cầm và hai cây gậy maraca gỗ, như thể đó là những bảo vật quý giá: ``Đây! Hai đứa kia sẽ chơi cái này ha!''

\img{6-5e}

Nàng Pháp sư trố mắt, bất bình phản ứng, giọng the thé: ``Bộ chị nghĩ chúng ta là một nhóm diễn hài chắc!?''

Cậu Tổng cười lớn, nói với tông xóc xỉa, giọng như một nghệ sĩ hài nhận xét về tiết mục mới: ``Thì mấy đứa lúc nào chả thế. Một gánh xiếc hề là đằng khác.''

Nghe vậy, nàng Pháp sư càng bực hơn, hai tay chống nạnh: ``\textbf{Anh đang nói cái quái gì thế-nanodesu!?}'', còn nàng Sói đen chỉ có thể vỗ vai an ủi như thể cô đã nghe cậu Tổng châm chọc như vậy nhiều lần như đã trở thành bản năng.

Nàng Lễ hội chỉ cười hì hì, vỗ vai cô đồng tóc xanh: ``Bọn chị toàn bị thế suốt mà.''

Nàng Sói đen thở dài, đôi tai cụp xuống: ``Dù gì anh ấy nói cũng\dots\ phần nào đúng mà\dots''

Cô cầm hai cây gậy maraca lên, xoay xoay trong tay rồi thử lắc nhẹ, tạo ra tiếng lạch xạch khô khốc. ``Cơ mà thứ này\dots'' cô nhíu mày, ``nhìn chả giống nhạc cụ tí nào cả. Cùng lắm nó chỉ có tiếng lạch xạch khi lắc thôi.''

Nàng Pháp sư khoanh tay gật gù: ``Nghe cứ như đồ chơi trẻ con ấy.''

Kuon đưa mắt nhìn chiếc kẻng mà nàng Lễ hội đang cầm trên tay, tò mò hỏi: ``Mà tiện thể thì, cái kẻng này ở đâu ra vậy?''

Matsuri vỗ ngực đầy tự hào, như thể vừa giành được giải thưởng lớn: ``AkiAki đưa cho em thứ này!''

Ryuuhou nghe vậy, khẽ nhướng mày, giọng tò mò: ``Aki Rosenthal-san - đồng môn của chị ấy có thú vui sưu tầm đồ cổ lạ, em nghe nói thế. Chắc đây là một món quà `thử thách' chăng?''

Kuon gật gù, đưa mắt nhìn em gái với vẻ thích thú: ``Em cũng biết Aki-chan cơ à?''

``Em có nghe qua vài lần từ mấy người bạn trong giới âm nhạc,'' – cô bé cười tinh quái, – ``chị ấy từng tặng một chiếc kèn đồng bị rỉ sét cho một ban nhạc khác. Kết quả là họ đổi tên thành `Đồng phế liệu' luôn.''

Kuon nheo mắt, giọng pha chút nghi ngờ, ngón tay gõ nhẹ lên thành ghế: ``\dots\ Nói thật, cái kẻng này thì làm được trò trống gì? Trông nó cũng chả khác một chiếc kẻng bình thường là mấy.''

Nàng Lễ hội cười rạng rỡ, hai mắt ánh lên vẻ bí ẩn như một phù thủy vừa tìm được bùa chú mới: ``Thì bởi vì nó có phải là cái kẻng bình thường đâu! Khi ta đánh vào cái kẻng này, tùy thuộc vào độ thuần khiết của tâm người đánh mà thanh của nó sẽ khác đó!''

\img{6-5f}

Mio nhướng mày, đôi tai dựng lên đầy thích thú: ``Ái chà.''

Rushia đặt tay lên cằm, tỏ vẻ hứng thú như một nhà nghiên cứu đang khám phá hiện tượng siêu nhiên: ``Một nhạc cụ nghe có vẻ thần kỳ nhỉ-nodesu.''

Ngồi ở cuối ghế, em gái của cậu Tổng khẽ nhấp một ngụm trà rồi lên tiếng, giọng mang theo chút hoài nghi chuyên môn: ``Theo kiến thức âm nhạc của em thì không có nhạc cụ nào phát ra âm thanh dựa trên `tâm hồn' cả. Chỉ có cách đánh và chất liệu thôi. Matsuri-san, chị có chắc đây không phải đồ chơi trẻ con chứ?''

Ông anh trai của cô gật gù, vẫn khoanh tay, tặc lưỡi tiếp lời em gái: ``Nói gì thì nói, nhưng nếu đánh nó cũng chỉ có một tiếng `keng' đơn điệu kêu lên thôi mà.''

Matsuri chu môi, không để tâm đến lời bình luận của hai anh em: ``Biết đâu được! Em đánh nha!''

Cô nàng Cổ động cầm dùi lên - trông còn mới hơn cả cái kẻng - vung tay gõ mạnh vào mặt kẻng với một khí thế dũng mãnh.

*BỤP.*

Một âm thanh trầm đục vang lên, như tiếng đập vào thùng giấy rỗng, rồi vụt tắt trong giây lát, nó tắt ngúm ngay trước khi kịp tạo ra bất kỳ cảm xúc nào.

Không có tiếng ngân, không có tiếng vang, cũng chẳng có chút gì gọi là ``thần kỳ''. Căn phòng rơi vào một sự im lặng đáng sợ, đến nỗi có thể nghe thấy tiếng lá vàng rơi ngoài sân.

Ba người còn lại trợn tròn mắt, không ai nói nổi câu nào. Cô nàng mắt hai màu cũng khẽ nhướng mày, nhưng rồi lại kìm nén nụ cười, chỉ lẳng lặng đặt tách trà xuống.

Matsuri cũng đờ người ra, miệng há hốc như vừa bị đánh gục bởi chính đòn của mình. Rồi chẳng nói chẳng rằng, cô kéo một tấm chăn từ đâu đó ra - như một ảo thuật gia với chiếc khăn trùm biểu diễn - quấn quanh người rồi lầm bầm, giọng như ma ám: ``\hl{Khi nào trời tối thì gọi tôi dậy nha\dots}''

Cô lập tức quay lưng, nằm cuộn tròn như con tôm trong vỏ chăn, trông hệt như một linh hồn bị tổn thương nặng nề vừa bị thế giới từ chối.

\img{6-5g}

Kuon khẽ nhếch mép, giọng châm chọc: ``Này Mat-chan\dots\ đừng nói là\dots'' Cậu không cần nói hết câu, nhưng cậu cũng đã đoán định được tại sao cô nàng lại như vậy: dường như tâm hồn của cô không được trong sáng cho lắm, trái với những gì cô tuyên bố.

Mio xóc hai cây gậy maraca lên như một công cụ tố cáo, hét lên đầy bất lực: ``\textbf{Em biết là chị đang bị sốc rồi, nhưng khoan ngủ đã!!}''

Cô quay sang cậu con trai, chỉ tay vào chiếc chăn như thể vừa phát hiện một bằng chứng phạm tội: ``Với lại, chị ấy kiếm đâu ra cái chăn vậy hả!?''

Cậu Tổng bình thản đáp, giọng điềm như nước hồ thu: ``Văn phòng mình có một bộ chăn gối sơ cua, đề phòng có ai đó ngủ lại mà. A-san dùng cái đó suốt.''

Nàng Sói đen nhướng mày, vẫn chưa hết kinh ngạc: ``Điều đó vẫn chẳng giải thích được tại sao chị ấy lại kiếm cái bộ đó nhanh đến thế!''

Lúc này, nàng Trưởng nhóm chen vào, giọng đầy vẻ trầm ngâm: ``Em nghĩ Matsuri-san có tài năng đặc biệt trong việc tìm chăn. Giống như em có tài tìm nhạc cụ rẻ vậy.''

Nhưng Matsuri không còn nghe thấy gì nữa, những lời bông đùa của Ryuuhou trôi qua như gió thoảng. Cô dỗi đến mức chẳng thèm phản ứng, chỉ vùi mặt vào chăn như thể cả thế giới này không còn quan trọng với mình nữa. Chỏm tóc ngố của cô rũ xuống, tràn trề nỗi thất vọng đến mức có thể làm khô cả một đại dương.

Mio bước đến gần Matsuri, vỗ nhẹ lên vai cô như một cách dỗ dành, giọng dịu dàng hơn bao giờ hết: ``Thôi nào, vui lên đi!''

``Đúng đó, cái này thì anh nghĩ chả việc gì mà phải dỗi cả,'' cậu Tổng lên tiếng, vẫn giữ giọng bình thản nhưng có phần cố an ủi.

Cô gái mắt hai màu gật nhẹ, bổ sung: ``Em đồng ý với anh ấy. Chị thử lần nữa xem, lần này đánh mạnh hơn, hoặc để em thử xem nào. Em đã từng đánh thử mấy cái kẻng trong phòng tập, âm thanh của nó phụ thuộc vào cách cầm dùi cơ.''

``Anh không nghĩ vấn đề là cách đánh của nó đâu, Ryuuhou\dots'' Kuon thở dài.

Mio kiễng người xuống, tiếp tục vỗ về nàng Lễ hội như một đứa trẻ: ``Nếu chị thực sự muốn chơi như thế, vậy chỉ cần thanh lọc tâm của chị là xong chứ gì!''

Nghe đến đây, Matsuri bỗng giật mình, tóc vểnh cả lên, ánh mắt ngờ vực nhìn Mio, như thể vừa nghe thấy một điều vô cùng mới lạ: ``Thanh lọc tâm của chị á?''

Cậu con trai gật đầu, giọng chắc nịch như một vị thầy đang ban phát chân lý: ``Bọn anh không biết sẽ làm thế nào, nhưng bọn anh sẽ cố gắng.''

``Đúng á! Bọn em sẽ giúp chị làm trong sạch nó!'' Mio đáp lời, đôi mắt sáng ngời quyết tâm.

Nàng Sói đen nói với vẻ chân thành, ánh mắt đầy nghị lực, như một chiến binh đang tuyên thệ trước trận đánh. Matsuri nhìn cô, trong mắt bắt đầu long lanh như thể sắp khóc, với giọt nước mắt sắp tràn ra vì xúc động.

\img{6-5h}

``Mio-chan à\dots''

Khung cảnh này trông không khác gì một người bạn đang động viên nhân vật chính khi họ thất bại cả: một màn kịch đầy cảm xúc giữa hai người bạn thân.

Rushia thở dài, đưa mắt nhìn Kuon với vẻ bất lực: ``Ta định nhìn hai chị ấy tình tứ với nhau đến bao giờ nữa?''

Cậu Tổng chỉ mỉm cười, bình thản đáp, như một đạo diễn đang tận hưởng cảnh quay: ``Cứ để họ yên đi.''

Em gái của cậu ngồi bên cạnh, ghé sát tai anh thì thầm: ``Em thấy đây là phim tình cảm tập dài đấy, Onii-san ạ. Có khi còn dài hơn cả buổi tập nhạc của bọn em.''

Cậu Tổng khẽ cười, đáp lại bằng ánh mắt tán đồng.

Sau một hồi hai người cứ nhìn nhau bằng ánh mắt sâu thẳm như trong phim truyền hình dài tập, cả hai bất chợt thay đổi thái độ, ánh mắt chuyển từ u sầu sang quyết tâm cháy bỏng, như hai ngọn lửa vừa được thổi bùng.

``\textbf{Ta cùng cố gắng nha!!}'' nàng Sói đen giơ tay lên trời đầy khí thế, giọng vang dội như tiếng kèn xung trận. Matsuri cũng lập tức hưởng ứng bằng cách giơ tay theo, chiếc chăn tuột khỏi vai, phất phơ như lá cờ chiến thắng.

Từ bên ngoài, Rushia, Kuon và Ryuuhou nhìn cảnh tượng này với những phản ứng khác nhau.

Nàng Pháp sư thì rùng mình, xoa trán, giọng đầy lo âu: ``Em có cảm giác bất an rồi đấy-nanodesu\dots''

Cậu Tổng thì chỉ nhún vai, khoanh tay đứng dựa vào tường: ``Cứ để xem họ sẽ làm gì\dots''

Cô em gái mỉm cười tinh quái, nhấp ngụm trà cuối cùng rồi đứng dậy: ``Nếu cần có người gõ kẻng đúng kỹ thuật, em sẵn sàng hỗ trợ đấy. Nhưng em không chịu trách nhiệm nếu nó lại phát ra tiếng `bụp' một lần nữa đâu.''

\ct

Một lúc sau, cả ba người khoanh tay, đứng từ xa quan sát cảnh tượng kỳ lạ trước mặt, như hai vị khán giả đang xem một vở kịch thiếu kịch bản.

Nàng Lễ hội đang ngồi khoanh chân trên sàn, nhắm mắt lại, hai tay đặt lên đầu gối, trông như một thiền sư tập sự vừa rời khỏi núi tuyết, thế nhưng chiếc mũ vẫn nằm lệch trên đầu khiến khung cảnh mất hết vẻ trang nghiêm. Nàng Sói đen đứng phía sau cô, hai tay đặt lên vai, miệng liên tục lẩm bẩm mấy câu thần chú kỳ quặc, giống hệt một pháp sư đang trừ tà trong phim kinh dị.

Kuon nhìn cảnh tượng trước mặt, nhíu mày, giọng đầy hoài nghi: ``Em có chắc là cách này sẽ được chứ?''

Mio gật đầu tự tin, đôi mắt sáng lên như đèn pha: ``Đảm bảo là sẽ được thôi, Kuon-senpai!''

Ngồi ở góc phòng, Ryuuhou vừa lướt điện thoại vừa khẽ nhướng mày: ``Thiền tịnh là tốt đấy, nhưng em phải nói thật là âm thanh của kẻng phụ thuộc vào chất liệu và lực đánh. Tâm hồn thanh tịnh mà gõ yếu thì vẫn ra tiếng bụp thôi.''

Rushia nghiêng đầu nhìn tiền bối của mình, không giấu nổi vẻ nghi ngại: ``Vì em thấy nãy giờ chị ấy chỉ ngồi thiền tĩnh tâm thôi\dots''

``Thì nó cũng hiệu quả mà,'' nàng Sói đen đáp, giọng đầy tự hào, như thể đã tìm ra chìa khóa của vũ trụ.

Nghe câu trả lời đó, cậu Tổng nuốt nước bọt, toát mồ hôi: ``Nhìn thế chứ anh thấy sợ vãi\dots''

Nàng Pháp sư gật đầu đồng ý, đôi tay siết chặt hơn. Cả hai người cùng đưa mắt nhìn nhau, như thể đang tự hỏi liệu có nên ngăn cản nghi lễ kỳ quặc này hay không, nhưng rồi ai cũng chọn cách im lặng vì tò mò hơn sợ hãi.

\ct

Sau một hồi tập thiền kéo dài đến mức có thể khiến một nhà sư cảm thấy tự hào, cuối cùng nàng Lễ hội đã có thể tự tin đánh chiếc kẻng thiêng này - ít nhất thì bản thân cô nghĩ vậy.

Chủ nhân của chiếc kẻng - Aki - xuất hiện từ cánh cửa, bước chân nhẹ nhàng như một linh hồn bay lướt. Cô tiến đến gần Matsuri, người đang cầm chiếc kẻng trong tay với tư thế như một pho tượng cổ, và hỏi với giọng nhẹ nhàng, đầy tò mò: ``Bà có chơi được không, Matsuri-chan?''

Cô nàng Lễ hội không trả lời mà chỉ nở một nụ cười thật tươi, mắt nhắm lại, như thể là một bức tượng phật trong trạng thái hỉ lạc vĩnh cửu; một nụ cười bất diệt đến mức có thể gây rùng mình cho bất kỳ ai nhìn vào.

Cậu Tổng nhìn khuôn mặt hỉ đó, thở dài bất lực, đưa tay xoa trán: ``Em ấy đã như thế khoảng 15 phút rồi.''

Nàng Sói đen, người đã tận mắt chứng kiến toàn bộ quá trình, cúi gằm mặt xuống tỏ vẻ mệt mỏi, đôi vai rũ xuống: ``Tâm của chị ấy trong sáng đến mức mỉm cười 24/7 luôn rồi\dots''

\img{6-5i}

Nàng Dị giới tròn mắt ngạc nhiên, hốt hoảng lui lại một bước: ``Ehh!? Sợ thế!''

Ngồi bên cửa sổ, cô nàng mắt hai màu nhìn Matsuri với vẻ thích thú, thì thầm với anh trai: ``Nhìn chị ấy như một vị phật cười trong chùa ấy. Có khi em nên xin chị ấy một chút `ánh sáng thiền' để áp dụng cho ban nhạc nhỉ.''

``Ít nhất thì cái tiếng kẻng chắc phải nghe hay hơn rồi nhỉ?'' cậu Tổng chống hông, nhìn vào chiếc kẻng trong tay Matsuri, giọng đầy hoài nghi.

Nàng Sói đen đứng dậy, đặt tay lên vai Matsuri, rồi hướng ánh mắt về chiếc kẻng trong tay cô nàng, giọng đầy hy vọng: ``Vậy em chơi thử đi!''

Nàng Dị giới hào hứng vỗ tay, như một đứa trẻ sắp được xem ảo thuật: ``Được rồi! Một, hai\dots''

Kuon, Aki và Mio bắt đầu đếm nhịp đồng thanh: ``Hai, ba!''

Matsuri với đôi mắt vẫn nhắm nghiền và nụ cười như pho tượng nâng chiếc gậy lên, rồi gõ mạnh một cái vào kẻng.

*KENGGGGGGGGGGGGGGGGGGGGGGGGGGGGGGGGGGG!! *

\gif{6-5j}{1}{152}

Lập tức, một thứ âm thanh vang vọng cả văn phòng, lan ra khỏi tòa nhà, len lỏi đến thành phố, và rồi cả Trái Đất, như một cơn sóng âm thanh cuốn phăng mọi thứ trên đường đi.

Âm thanh trong trẻo và vang vọng đến mức người viết câu chuyện này ở Hà Nội còn có thể nghe thấy (!?) một minh chứng sống động cho sức mạnh vượt không gian, thời gian của âm thanh.

Cô em gái của cậu Tổng vội bịt tai, nhưng không giấu nổi vẻ kinh ngạc: ``Ôi trời, tiếng này vượt quá tần số của kẻng thường. Rõ ràng nó đang cộng hưởng với một cái gì đó\dots\ nhưng mà là gì nhỉ?''

\ct

Khoảnh khắc đó, nàng Aki và Mio cũng đột nhiên mỉm cười y như Matsuri, một nụ cười vô thức nhưng lại mang theo sự thanh thản lạ thường, như thể vừa chạm vào cõi niết bàn.

\img{6-5k}

Trong khi đó, Rushia và Kuon đứng chôn chân tại chỗ, trố mắt nhìn ba người bọn họ với sự ngỡ ngàng và khó hiểu, như thể vừa chứng kiến một phép thuật ngoài sức tưởng tượng.

Nàng Dị giới mỉm cười rạng rỡ, tuyên bố với một sự phấn khích đến khó tin, giọng như tiếng chuông báo tin vui: ``Và với tiếng kẻng đó, bọn chị đã lập nhóm `Chị em Hút bụi Ion âm' hoàn chỉnh!''

Nàng Pháp sư há hốc miệng, trợn mắt: ``Vậy cũng được nữa!?'' Cô hoàn toàn không thể tin vào những gì mình vừa nghe, tay xoa xoa thái dương như thể đang cố gắng tiêu hóa mớ thông tin phi logic này.

Cậu Tổng thở dài, đưa tay lên trán như thể đang cố gắng tiêu hóa toàn bộ sự việc hoang đường này, giọng đầy bất lực: ``Anh không nghĩ là nó có tác dụng đấy\dots\ Với cả mấy đứa vẫn nhất quyết dùng cái tên đó sao?''

Ngồi bên cạnh, Ryuuhou nhún vai, giọng đầy thích thú: ``Em phải công nhận là cái tên có sức hút riêng. Nếu đổi thành `Ion Negative Girls' thì nghe sang hơn, nhưng thôi, đó là ý kiến chuyên môn của em.''

Cả căn phòng rơi vào một bầu không khí kỳ lạ, khi một nửa số người chìm trong sự thanh tịnh đến mức tưởng chừng có thể bay bổng, còn nửa còn lại thì ngập trong hoang mang như lạc giữa một cơn ác mộng màu hồng.

``\textit{Cơ mà mình thật sự chả biết đây là nhóm nhạc hay là nhóm tấu hài gì nữa\dots}'' Kuon lẩm bẩm, lắc đầu, nhưng đã bắt đầu thấy tò mò về màn tiếp theo của họ.

%==========================================TRUYỆN PHỤ=========================================
\omk

Với đội hình ba người: Aki chơi vĩ cầm, Mio cầm bộ maraca và Matsuri cầm kẻng, họ đã lập một ban nhạc\dots\ mà cũng chẳng ra ban nhạc cho lắm – nó giống một tiết mục xiếc hơn là một buổi hòa nhạc.

Rushia thở dài, nhìn cả ba người đang nở một nụ cười thanh tao kia, giọng e sợ như thể đang nhìn vào một bức tượng có sức mạnh siêu nhiên: ``Liệu với nhóm này thì họ có thể làm được trò trống gì không nữa.''

Kuon nhăn mặt, gật đầu đồng tình, đưa tay xoa cằm: ``Anh cũng nghi lắm\dots''

Cả hai người ngoài cuộc đứng nhìn ba con người trong nhóm nhạc mới ``ra đời'', cảm giác như thể đang quan sát một vở hài kịch kỳ quái vừa ra mắt khán giả. Nàng Lễ hội đang cười tươi rói, đôi mắt nhắm tịt và gương mặt như pho tượng cười vĩnh cửu, nàng Dị giới cũng thả lỏng và miệng không ngừng hát theo tiếng kẻng đều đều của Matsuri – một giai điệu ``keng\dots\ keng\dots'' lặp đi lặp lại không có hồi kết, trong khi nàng Sói đen tung hứng những chiếc maraca lên không trung như thể đang tham gia một buổi biểu diễn xiếc lễ hội, quay cuồng như một diễn viên múa rối.

Ngồi nép ở ghế sô-pha, Ryuuhou cất giọng bình luận, hai tay chống cằm: ``Tiết tấu của Matsuri-san bị lệch nhịp khoảng nửa phách so với tiếng maraca của Mio-san, còn cây vĩ cầm của Aki-san thì lạc sang tông khác. Nếu đây là một buổi tập của ban nhạc em, em đã cho họ về nhà ngay lập tức rồi.''

Cậu Tổng thở dài, than thở, giọng đầy bất lực như một vị đạo diễn đang chứng kiến một thảm họa phim ảnh: ``\dots Anh lại thấy trông họ như những diễn viên hài kịch hơn là nhạc công đấy.''

Nàng Pháp sư lắc đầu, tay ôm trán: ``Matsuri-senpai và Aki-senpai thì em còn hiểu, nhưng Mio-senpai, vốn dĩ nghiêm túc thế cũng hùa vào mới hay chứ\dots''

Chàng quản lý gật gù, ngán ngẩm nhìn nàng Sói đen vốn đã từng nghiêm túc trước kia – giờ đây đang tung hứng gậy như một chú hề vui vẻ: ``Chiều lòng em ấy là chính, chứ thực sự anh cũng nghĩ là Mio-chan cũng không muốn làm đâu. Nhưng nhìn họ vui thế kia là anh cũng đành chịu\dots''

Cô em gái khoanh tay, nhìn với ánh mắt phân tích: ``Theo em, tâm lý đám đông đã thôi thúc chị ấy tham gia. Ở trong một ban nhạc, sự hào hứng của người khác dễ lây lan, cả cái tích cực lẫn tiêu cực. Có lẽ chị Mio đang trong giai đoạn `nhiễm bệnh vui vẻ' đấy.''

Cả ba người trong nhóm bắt đầu làm chơi thứ nhạc cụ mà họ cầm trên tay, nhưng tiếng động phát ra lại chẳng hề giống âm nhạc; nó giống một bản giao hưởng của một đám trẻ con trong giờ thể dục hơn. Nàng Sói đen quay chiếc lục lạc thành vòng tròn trên tay, khiến nó phát ra tiếng ``lạch xạch'' kỳ lạ mà không hề tạo ra một giai điệu nào, như thể một con sói đang cố gắng nhảy múa mà không biết nhạc. Nàng Dị giới, với vĩ cầm trong tay, cũng không có gì khả quan hơn, khi cô nàng chỉ gảy thử vài dây với nụ cười thánh thiện mà không tạo ra được một bản nhạc hoàn chỉnh, mà chỉ là những nốt đơn lẻ ngẫu nhiên như tiếng mèo kêu. Còn nàng Lễ hội, với chiếc kẻng, thì chẳng khác gì một gánh xiếc đang biểu diễn trong buổi tổng duyệt cuối cùng trước khi bị hủy.

Matsuri hứng thú, thốt lên vui vẻ, giọng đầy tự tin như thể mình vừa tìm ra một quy luật mới của vũ trụ: ``Xem đây, bọn mình sẽ tạo ra một bản giao hưởng siêu phàm này!''

Tuy nhiên, tiếng ``keng'' do cô tạo ra giờ đây lại chẳng có gì đặc biệt, chỉ làm cho không khí thêm phần hỗn loạn, như một tiếng chuông báo động giữa một buổi hòa nhạc đang tồi tệ. Cuối cùng, một giọng hát của nàng Pháp sư bất chợt vang lên. Cô đồng không thể chịu đựng được nữa - bị kéo vào cơn cuồng loạn âm thanh đó - đã quyết định sẽ tham gia và hát, dù cô không hiểu gì về việc chơi nhạc, chỉ đơn thuần hét lên để giải tỏa sự bực bội.

``\textbf{Tại sao tôi lại phải làm theo những cái trò này chứ!-nanodesu?}''

Dù vậy, cô nàng vẫn phải gượng gạo bắt đầu cất giọng, một giọng hát giữa lời than vãn và tiếng la hét, nhưng điều đó lại hòa hợp một cách kỳ lạ với phần nền hỗn loạn. Và kết quả là, một buổi trình diễn điên loạn bắt đầu, một tác phẩm nghệ thuật ngẫu nhiên.

\img{6-5l}

Mio cứ tiếp tục tung hứng chiếc gậy gỗ không ngừng, tiếng ``lạch xạch'' cứ vang lên như một đoạn nhạc điện tử hỏng hóc trong khi Matsuri gõ vào chiếc kẻng với nhịp điệu không đâu vào đâu - một cú sốc cho bất kỳ ai có năng khiếu âm nhạc. Akicũng không chịu thua, tiếp tục gảy vĩ cầm một cách bất chấp, như thể cô đang luyện tập cho một buổi biểu diễn mà không ai mời. Và đương nhiên, không thể thiếu được tiếng hét của Rushia, thứ âm thanh giống như một tiếng la hét chói tai, khiến cả văn phòng trở thành một hỗn loạn tồi tệ, tựa như một bản giao hưởng của sự điên rồ.

Cả bốn người không hề nhận ra rằng âm thanh họ phát ra có thể phá vỡ yên tĩnh của cả một khu vực, thậm chí làm rung cả cửa sổ, mà có lẽ một vài người hàng xóm đã gọi cảnh sát từ lúc nào chẳng ai biết. Bộ maraca của nàng Sói đen bị ném lên trời rồi rơi xuống đất với một tiếng ``maraca~{}'' vang lên liên tục, khiến không gian xung quanh như bị xé toạc. Nàng Lễ hội gõ kẻng như thể đó là một phần của một màn trình diễn trong lễ hội: mạnh mẽ và đều đặn, nhưng chẳng có một chút nghệ thuật nào.

``\textbf{Coi chừng ta đây! Ta đây là Cô đồng triệu hồi từ cái chết!!}'' nàng Pháp sư hét lớn, dường như đã hòa vào bản ``hùng ca'' của sự lộn xộn, giọng cô vang lên như một lời nguyền rủa trong một bộ phim kinh dị.

Nàng Dị giới thì vui vẻ tiếp tục gảy những đoạn violin lạc quẻ và lệch nhịp, nụ cười thánh thiện của cô như đang ban phước cho sự hỗn loạn: ``Cứ thoải mái đi mà!''

Ngồi một góc khác, nàng Trưởng nhóm cười khúc khích, lắc đầu với vẻ hài lòng: ``Âm nhạc của họ không có cấu trúc, không có tông, không có tiết tấu\dots\ và cũng không có khán giả. Nhưng thành thật mà nói, em vẫn thấy thú vị hơn mấy buổi tập nhàm chán của ban nhạc em đấy.''

Trong khi đó, cậu Tổng đứng ở góc phòng, mặt mày hốc hác như vừa nhìn thấy ngày tận thế, bịt chặt tai nghe trong tai để cố gắng không nghe thấy cái tiếng hỗn loạn ấy nhưng vô ích, âm thanh len lỏi qua mọi kẽ hở. Cậu cố gắng tập trung vào công việc của mình, nhưng âm thanh xung quanh cứ không ngừng quấy rầy, mỗi tiếng ``keng'', ``lạch xạch'', ``re re'' lại đập vào màng nhĩ như một lời nhắc nhở về sự điên rồ.

Cậu nhắm mắt, thở dài, giọng như một lời cầu nguyện vô vọng: ``Họ nên có một buổi học nhạc cơ bản trước đi đã\dots''

Cậu thở dài trong ngao ngán, cảm giác như mình đã bị cuốn vào một tình huống không thể thoát ra, như thể đang lạc trong một cơn ác mộng màu sắc sặc sỡ. Tất cả những gì chàng quản lý có thể làm là ngồi im, không dám nhìn về phía những người đang tạo ra cái hỗn loạn đó, sợ rằng một cái nhìn sẽ khiến mình bị hút vào vòng xoáy không lối ra.

Em gái của cậu cuối cùng cũng không thể kìm được, cô đứng dậy, cầm một cây đàn ukulele nhỏ xíu từ trong cặp ra và bắt đầu đệm vào với một đoạn riff ba hợp âm đơn giản nhưng đúng nhịp, hòa vào bản nhạc hỗn loạn như một điểm sáng nhỏ giữa bão tố.

Kuon nhìn cô em gái, không khỏi ngạc nhiên: ``Em cũng nhập hội luôn hả, Ryuuhou!?''

Cô nàng mắt hai màu cười tinh ranh, giọng đầy hứng khởi: ``Điên rồ chung vui mà, Onii-san! Hơn nữa em thấy họ đang cần một ai đó có chút kỹ thuật để giữ họ khỏi lao xuống vực, và trông này, em đang làm được đấy!''

Và thế là, một buổi sáng làm việc bình yên đã bị phá vỡ hoàn toàn, nhưng bù lại, những tiếng cười và sự vui vẻ ấy lại làm cho văn phòng trở nên đầy ắp một bầu không khí đặc biệt\dots\ dù có chút hỗn loạn và kỳ quặc:  một bản giao hưởng của tình bạn, sự ngớ ngẩn và niềm vui bất tận.

\end{document}